\documentclass[12pt,reqno]{amsart}
\usepackage{math-review}
% Project-wide mathematical notation. No claims are certified here.

\title{Internal Paper-Mode Boundary Example}
\date{}
\begin{document}\maketitle
\noindent\textbf{DRAFT\_UNVERIFIED: not an original research paper; not for formal publication.}
\begin{abstract}This document exercises the research-paper software boundary using known
linear algebra. No original theorem, prior-art assessment or submission-ready research
quality is claimed. A complete proof of a classical statement cannot supply originality.\end{abstract}
\section{Known result, not a contribution}
\begin{theorem}[Classical criterion]
Let $p,q\geq1$ be integers. Let $A\in\mathbb R^{p\times p}$ and
$C\in\mathbb R^{q\times q}$ be symmetric, and let $B\in\mathbb R^{p\times q}$.
Assume that $x^{\mathsf T}Ax>0$ for every $0\ne x\in\mathbb R^p$. Define
\[
 M=\begin{pmatrix}A&B\\ B^{\mathsf T}&C\end{pmatrix},
 \qquad S=C-B^{\mathsf T}A^{-1}B.
\]
Then $M$ is positive definite if and only if $S$ is positive definite.
Here a real symmetric matrix $H$ is positive definite when
$v^{\mathsf T}Hv>0$ for every nonzero real vector $v$ of the appropriate size.
\end{theorem}
\section{Established inputs}
\begin{lemma}Let $A\in\mathbb R^{p\times p}$ be symmetric and invertible,
$B\in\mathbb R^{p\times q}$ and $C\in\mathbb R^{q\times q}$ symmetric.
For all $x\in\mathbb R^p$ and $y\in\mathbb R^q$, putting
$z=x+A^{-1}By$ gives
\[
 x^{\mathsf T}Ax+2x^{\mathsf T}By+y^{\mathsf T}Cy
 =z^{\mathsf T}Az+y^{\mathsf T}(C-B^{\mathsf T}A^{-1}B)y.
\]\end{lemma}
\begin{lemma}For every real matrix $K\in\mathbb R^{p\times q}$, the linear map
\[
 (x,y)\longmapsto (x+Ky,y)
\]
is invertible, with inverse $(z,y)\mapsto(z-Ky,y)$. Consequently it takes
nonzero vectors to nonzero vectors, in both directions.\end{lemma}
\begin{proposition}Let $A\in\mathbb R^{p\times p}$ be symmetric and invertible,
$B\in\mathbb R^{p\times q}$ and $C\in\mathbb R^{q\times q}$ symmetric.
For
\[
 M=\begin{pmatrix}A&B\\B^{\mathsf T}&C\end{pmatrix},\quad
 T=\begin{pmatrix}I_p&-A^{-1}B\\0&I_q\end{pmatrix},\quad
 S=C-B^{\mathsf T}A^{-1}B,
\]
one has $T^{\mathsf T}MT=\operatorname{diag}(A,S)$.
Because $T$ is invertible, $M>0$ is equivalent to $A>0$ and $S>0$ jointly;
$S>0$ alone is not sufficient without $A>0$.\end{proposition}
\section{Proof and publication boundary}
Write $Q(x,y)=x^{\mathsf T}Ax+2x^{\mathsf T}By+y^{\mathsf T}Cy$.

For $A,S>0$, the square identity and invertible coordinate change show that
$Q(x,y)>0$ for every nonzero pair. Conversely, for every nonzero $y$, set
$x=-A^{-1}By$; then $M>0$ implies $y^{\mathsf T}Sy>0$.
The mathematical argument is closed but known. The missing original contribution
and prior-art assessment are disclosed, not filled with fictional novelty records.
\nocite{S-guide}\bibliographystyle{plain}\bibliography{../generated/references}
\end{document}
